% =========================================================================
% ภาคผนวก จ: คู่มือขั้นตอนการปฏิบัติงานมาตรฐานภาคสนาม (Field Protocols & SOPs)
% Appendix E: Field Protocols and Standard Operating Procedures (SOPs)
% =========================================================================

\chapter{คู่มือขั้นตอนการปฏิบัติงานมาตรฐานภาคสนาม (Field Protocols \& SOPs)}
\label{app:field_protocols}

ภาคผนวกนี้จัดทำขึ้นเพื่อเป็นคู่มือขั้นตอนการปฏิบัติงานมาตรฐาน (Standard Operating Procedures, SOPs) สำหรับนักวิจัย เกษตรกรปราดเปรื่อง (Smart Farmers) และผู้ประกอบการนวัตกรรมเกษตร โดยระบุขั้นตอนการทำงานอย่างเป็นลำดับขั้นตอน การควบคุมตัวแปรทางวิศวกรรม อุปกรณ์ที่ต้องใช้ และเกณฑ์การตรวจสอบคุณภาพผลผลิต เพื่อให้สามารถถ่ายทอดเทคโนโลยีสู่การปฏิบัติจริงได้อย่างถูกต้องและปลอดภัย

\section{SOP-01: ขั้นตอนการผลิตและสอบเทียบน้ำกระตุ้นด้วยพลาสมา (PAW Preparation)}
\label{app:sop_paw}

\index{ระเบียบวิธีปฏิบัติงานมาตรฐาน}\index{น้ำกระตุ้นด้วยพลาสมา}\index{การสอบเทียบเซนเซอร์}

\begin{enumerate}[leftmargin=2.0em, itemsep=6pt]
    \item \textbf{วัตถุประสงค์:} เพื่อผลิตน้ำกระตุ้นด้วยพลาสมาที่มีความเข้มข้นของไอออนไนเตรต ($\ce{NO3^-} \ge 200\ \si{\milli\gram\per\liter}$) ไฮโดรเจนเปอร์ออกไซด์ ($\ce{H2O2} \ge 50\ \si{\milli\gram\per\liter}$) และค่าศักย์ออกซิเดชันรีดักชัน ($\text{ORP} \ge +550\ \si{\milli\volt}$)
    \item \textbf{อุปกรณ์และสารเคมีที่ต้องใช้:}
    \begin{itemize}
        \item เครื่องกำเนิดพลาสมาเจ็ทความดันบรรยากาศ (CAP Jet) หรือระบบ Gliding Arc Discharge
        \item น้ำปราศจากไอออน (DI Water) หรือน้ำบาดาลกรองตะกอนความกระด้างต่ำ ($< 50\text{ ppm CaCO}_3$)
        \item เครื่องวัดมัลติพารามิเตอร์ภาคสนามสำหรับวัด pH, ORP, และ Electrical Conductivity (EC)
        \item ชุดทดสอบแผ่นเทียบสีหรือสเปกโตรโฟโตมิเตอร์วัดปริมาณ $\ce{NO3^-}, \ce{NO2^-}, \ce{H2O2}$
    \end{itemize}
    \item \textbf{ขั้นตอนการปฏิบัติงาน:}
    \begin{enumerate}[label=\arabic*., leftmargin=1.5em, itemsep=3pt]
        \item ล้างถังทำปฏิกิริยาด้วยน้ำ DI 2 ครั้งเพื่อกำจัดสิ่งปนเปื้อนอินทรียวัตถุตกค้าง
        \item เติมน้ำสะอาดปริมาตร $5.0\text{ ลิตร}$ ลงในถังปฏิกรณ์พลาสมา
        \item เปิดปั๊มหมุนเวียนสารละลาย (Peristaltic Pump) ที่อัตราการไหล $1.2\ \si{\liter\per\minute}$
        \item ปรับอัตราการไหลของอากาศอัดเข้าสู่หัวพลาสมาที่ $15.0\ \si{\liter\per\minute}$
        \item เปิดสวิตช์จ่ายกำลังไฟฟ้าแรงสูง ปรับตั้งค่า $V_{\text{peak}} = 8.5\text{ kV}$, ความถี่ $f = 25.0\text{ kHz}$
        \item ดำเนินการยิงพลาสมาต่อเนื่องเป็นเวลา $20\text{ นาที}$
        \item สุ่มเก็บตัวอย่างน้ำ $50\ \si{\milli\liter}$ วัดค่า pH, ORP ทันที และบันทึกผลลงในสมุดบันทึกแล็บ
    \end{enumerate}
    \item \textbf{เกณฑ์การยอมรับคุณภาพ (Acceptance Criteria):} ค่า $\text{pH}$ อยู่ในช่วง $2.80$--$3.40$, $\text{ORP} \ge +550\text{ mV}$, ค่าสภาพนำไฟฟ้า $\text{EC} \ge 1200\ \si{\micro\siemens\per\centi\meter}$ หากไม่ผ่านเกณฑ์ ให้ตรวจสอบการอุดตันของหัวพ่นอากาศหรือการเสื่อมสภาพของขั้วไฟฟ้า
\end{enumerate}

\section{SOP-02: ขั้นตอนการปรับสภาพเมล็ดพันธุ์ด้วยพลาสมา (Plasma Seed Priming)}
\label{app:sop_seed}

\index{การปรับสภาพเมล็ดพันธุ์}\index{ความงอก}\index{ความชอบน้ำ}

\begin{enumerate}[leftmargin=2.0em, itemsep=6pt]
    \item \textbf{วัตถุประสงค์:} เพื่อเพิ่มอัตราความงอก (Germination Rate) ดัชนีความแข็งแรง (Seed Vigor Index) และกระตุ้นการสร้างเอนไซม์ต้านอนุมูลอิสระในเมล็ดพันธุ์
    \item \textbf{ขั้นตอนการปฏิบัติงาน:}
    \begin{enumerate}[label=\arabic*., leftmargin=1.5em, itemsep=3pt]
        \item คัดเลือกเมล็ดพันธุ์ที่มีความสมบูรณ์ ปราศจากรอยแตกร้าวและความชื้นเริ่มต้น $10$--$12\%$
        \item กระจายเมล็ดพันธุ์บนถาดรองแก้วไดอิเล็กทริกเป็นชั้นบางระนาบเดี่ยว (Monolayer) ห้ามซ้อนทับกัน
        \item นำถาดเข้าสู่เครื่องปฏิกรณ์ Dielectric Barrier Discharge (DBD)
        \item ตั้งค่าระยะห่างช่องว่างการคายประจุ (Discharge Gap) เท่ากับ $3.0\ \si{\milli\meter}$
        \item จ่ายแรงดันพลาสมาความดันบรรยากาศที่ความหนาแน่นกำลังงาน $0.85\ \si{\watt\per\centi\meter\squared}$
        \item ทำการปรับสภาพเป็นเวลา $60\text{ วินาที}$ (สำหรับเมล็ดพืชเปลือกบาง) หรือ $180\text{ วินาที}$ (สำหรับเมล็ดเปลือกหนา)
        \item นำเมล็ดออกจากเครื่อง วัดค่ามุมสัมผัสหยดน้ำ (Water Contact Angle) ต้องลดลงเหลือ $< 45^\circ$
    \end{enumerate}
\end{enumerate}

\section{SOP-03: การตรวจวินิจฉัยและสอบเทียบสเปกตรัม OES ภาคสนาม}
\label{app:sop_oes}

\index{สเปกโทรสโกปีการปลดปล่อยเชิงแสง}\index{การสอบเทียบสเปกตรัม}

\begin{enumerate}[leftmargin=2.0em, itemsep=6pt]
    \item \textbf{วัตถุประสงค์:} เพื่อตรวจสอบอุณหภูมิการหมุน ($T_{\text{rot}}$) และความเข้มของการสลายตัวของโมเลกุลไนโตรเจนที่ความยาวคลื่น $337.1\text{ nm}$ แบบเวลาจริง (Real-time)
    \item \textbf{ขั้นตอนการสอบเทียบ:}
    \begin{enumerate}[label=\arabic*., leftmargin=1.5em, itemsep=3pt]
        \item ติดตั้งหัวรับสัญญาณใยแก้วนำแสง (Optical Fiber Probe) ทำมุม $90^\circ$ กับลำพลาสมา ห่างจากหัวฉีด $5.0\ \si{\milli\meter}$
        \item ปรับแต่งศูนย์สเปกโตรมิเตอร์ด้วยหลอดกำเนิดแสงสเปกตรัมปรอท-อาร์กอน (Hg-Ar Calibration Lamp)
        \item ตรวจวัดสัญญาณมืด (Dark Current Spectrum) และทำการหักลบอัตโนมัติในซอฟต์แวร์ประมวลผล
        \item ตรวจสอบความเข้มของเส้นสเปกตรัม $\ce{N2}(C^3\Pi_u \to B^3\Pi_g, \Delta v = 0)$ ที่ $337.1\text{ nm}$ หากพบว่าความเข้มลดลงมากกว่า $20\%$ ให้ทำความสะอาดช่องหน้าต่างควอตซ์
    \end{enumerate}
\end{enumerate}

\section{SOP-04: ข้อปฏิบัติการพ่นหมอก PAW ควบคุมเชื้อราในแปลงทุเรียน}
\label{app:sop_durian}

\index{ทุเรียน}\index{ไฟทอปธอรา}\index{การควบคุมทางชีวภาพ}

\begin{enumerate}[leftmargin=2.0em, itemsep=6pt]
    \item \textbf{การเตรียมสารละลายสำหรับฉีดพ่น:} ใช้น้ำ PAW สดใหม่ที่ผลิตไม่เกิน $6\text{ ชั่วโมง}$ โดยเจือจางด้วยน้ำสะอาดในอัตราส่วน $1:1$ สำหรับการพ่นใบ และใช้สารละลายเข้มข้น $100\%$ สำหรับการราดโคนต้น
    \item \textbf{ช่วงเวลาการพ่น:} ทำการฉีดพ่นในช่วงเช้าตรู่ ($06.00$--$08.30\text{ น.}$) หรือช่วงเย็นแดดร่ม เพื่อป้องกันการสลายตัวของสาร $\ce{ONOOH}$ และ $\ce{H2O2}$ จากรังสีความร้อนจากดวงอาทิตย์
    \item \textbf{การประเมินประสิทธิภาพ:} ติดตามตรวจสอบแผลเน่าบนเปลือกลำต้นทุเรียนหลังการพ่นทุก 7 วัน แผลเน่าแห้งและตกสะเก็ดสีน้ำตาลเข้มแสดงถึงการยับยั้งการเจริญของเส้นใยเชื้อรา *Phytophthora palmivora* ได้อย่างสมบูรณ์
\end{enumerate}

\clearpage
